\documentclass[12pt,british,round,comma,sort&compress]{article}
\renewcommand{\familydefault}{\rmdefault}
\usepackage[a4paper]{geometry}
\usepackage{graphicx}
\geometry{verbose,tmargin=2cm,bmargin=2cm,lmargin=2cm,rmargin=2cm}
\setlength{\parskip}{\smallskipamount}
\setlength{\parindent}{0pt}

\usepackage{color}
\newcommand{\nicefrac}[2]{\ensuremath ^{#1}\!\!/\!_{#2}}
\usepackage { fancybox}

\renewcommand{\floatpagefraction}{0.95}
\renewcommand{\textfraction}{0}
\renewcommand{\topfraction}{1}
\renewcommand{\bottomfraction}{1}

\usepackage{bibentry}
\bibliographystyle{abbrvnat}
\nobibliography{numerics}

\begin{document}

\title{Literature review relevant for Splitting}
\author{Hilary Weller}
\maketitle

\begin{itemize}

\item \bibentry{TV12}

\item \bibentry{YSSK24}
\item\bibentry{SHG90}
\item\bibentry{SYG90}
\item\bibentry{ZA20}
\item\bibentry{Thu24}

\item Email from Amber
\begin{quote}
Hilary and I thought it was maybe nice for you to have a reference for the GMRES(m) and GCR(k) algorithms. These are iterative matrix solvers to solve the matrix A [matrix equation Ax=b] connected to the implicit part of the time stepping, I don’t know how much you and Hilary have talked about this already - we can talk about it on Thursday if you like!

The key reference is Saad and Schultz 1986, which is about GMRES and GMRES(m). They also discuss GCR, though not GCR(k) directly. \url{https://epubs.siam.org/doi/10.1137/0907058}
The second key reference is Elman’s PhD thesis, which discusses GCR and GCR(k) (p. 32-34). \url{http://www.cvc.yale.edu/publications/techreports/tr229.pdf}
There are also numerous references online, iterative methods are widely used.

I use them to solve a heptadiagonal matrix for a fifth-order accurate spatial discretisation. Both work, though GCR may break down, because it requires a matrix with a positive definite symmetric part (SS1986) (I have not actually found any problematic cases myself). GMRES is also more efficient than GCR (SS1986).

FYI 1 : GCR and GCR(k) are pretty much the same thing, but the (k) means that the algorithm is restarted k times. Similarly for GMRES and GMRES(m): m is the number of restarts. The key reason for restarting is that it reduces the storage and number of calculations required.

FYI 2 : GCR stands for ‘Generalised Conjugate Residual’ and GMRES stands for ‘Generalised Minimal RESidual’.

FYI 3 : These are all Krylov methods, which e.g. ‘conjugate gradient’ and ‘conjugate residual’ are part of as well.

I can provide code that implements GMRES(m) and GCR(k) if you would like! But I’ll leave it be for now
\end{quote}

%\item \bibentry{} \ \\
%\item \bibentry{} \ \\
%\item \bibentry{} \ \\
%\item \bibentry{} \ \\

\end{itemize}

\nobibliography{numerics}

\end{document}
